% ---
% PACOTES BÁSICOS
% ---
\usepackage{lmodern}			% Usa a fonte Latin Modern
\usepackage[T1]{fontenc}		% Seleção de códigos de fonte
\usepackage[utf8]{inputenc}		% Codificação do documento (conversão automática dos acentos)
\usepackage{indentfirst}		% Indenta o primeiro parágrafo de cada seção
\usepackage{color}				% Controle das cores
\usepackage{graphicx}			% Inclusão de gráficos
\usepackage{microtype} 			% para melhorias de justificação
\usepackage{lipsum}				% para geração de dummy text
\usepackage{amsmath}			% para equações matemáticas
\usepackage{booktabs}			% para tabelas com padrão profissional
\usepackage{listings}			% para inserção de código-fonte
\usepackage{float}              % para forçar a posição de figuras (H)
\usepackage{caption}            % para legendas
\usepackage{csquotes}           % necessário para o biblatex

% ----
% Configuração de Quadros
\usepackage{newfloat}
% ---
% CONFIGURAÇÕES DO PACOTE DE CÓDIGO FONTE (listings)
% ---
\lstset{
  basicstyle=\ttfamily\footnotesize,
  breaklines=true,
  frame=single,
  numbers=left,
  numberstyle=\tiny,
  captionpos=b,
  keepspaces=true
}

% ---
% PACOTES DE CITAÇÕES E REFERÊNCIAS (BIBLATEX-ABNT)
% ---
% Utilizando biblatex-abnt, a solução mais moderna e compatível com Overleaf para normas ABNT
\usepackage[
  backend=biber,
  style=abnt,
  justify=true,
  sorting=nty
]{biblatex}
\addbibresource{referencias.bib}

% ---
% CONFIGURAÇÕES DE APARÊNCIA DO PDF FINAL
% ---
\definecolor{blue}{RGB}{41,5,195}

\makeatletter
\hypersetup{
     	%pagebackref=true,
		pdftitle={\@title}, 
		pdfauthor={\@author},
    	pdfsubject={\imprimirpreambulo},
	    pdfcreator={LaTeX with abnTeX2},
		pdfkeywords={abnt}{latex}{abntex}{abntex2}{relatório técnico}, 
		colorlinks=true,       		% false: boxed links; true: colored links
    	linkcolor=black,          	% color of internal links
    	citecolor=black,        		% color of links to bibliography
    	filecolor=magenta,      		% color of file links
		urlcolor=blue,
		bookmarksdepth=4
}
\makeatother

% ---
% ESPAÇAMENTOS
% ---
% O tamanho do parágrafo é dado por:
\setlength{\parindent}{1.25cm}

% Controle do espaçamento entre um parágrafo e outro:
\setlength{\parskip}{0.0cm}

% Compila o indice
\makeindex

% ---
% COMANDO \source (Fonte das ilustrações e tabelas)
% ---
\newcommand{\source}[1]{\vspace{0.2cm}\par\noindent\textbf{Fonte:} #1}


%----------------------------------------------------
% AMBIENTE QUADRO
%----------------------------------------------------

\DeclareFloatingEnvironment[
fileext=loq,
listname={Lista de Quadros},
name=Quadro,
placement=htbp
]{quadro}

% legenda acima do quadro
\captionsetup[quadro]{position=top}

% centraliza a legenda
\captionsetup[quadro]{justification=centering}

% fonte da legenda
\captionsetup[quadro]{font=normalfont}

% rótulo em negrito
\captionsetup[quadro]{labelfont=bf}

%----------------------------------------------------
% LISTA DE QUADROS
%----------------------------------------------------

\newcommand{\listadequadros}{
    \pdfbookmark[0]{\listquadroname}{loq}
    \listof{quadro}{\listquadroname}
}

\newcommand{\listquadroname}{Lista de Quadros}
